\chapter{Diazepam}

\section*{Definition and Overview}
Diazepam is a medicine from the benzodiazepine family that acts on the brain to reduce anxiety, relax muscles, stop seizures, and help with alcohol withdrawal. It is also used before some medical procedures to reduce anxiety, and in emergency situations such as prolonged seizures. Diazepam works by enhancing the effect of a natural brain chemical that calms the nervous system. It is a controlled medicine in Australia because it can cause dependence and misuse. Diazepam is effective for short-term use but is generally not recommended for long-term treatment of anxiety. This chapter explains how diazepam works, its uses, and its risks.

\section*{Epidemiology}
Benzodiazepines such as diazepam are among the most prescribed medicines in Australia, although prescribing has declined with greater awareness of dependence risks. They are used for anxiety, insomnia, muscle spasm, seizures, and alcohol withdrawal. Concern about dependence, falls in older people, and misuse has led to stricter prescribing and PBS restrictions. Diazepam remains an important emergency medicine for seizures and is carried by many people with epilepsy.

\section*{Aetiology and Pathophysiology}
\begin{itemize}
    \item \textbf{Mechanism:} diazepam enhances the action of GABA, the brain's main inhibitory neurotransmitter, calming nerve activity
    \item \textbf{Anxiolytic effect:} reduces anxiety and agitation
    \item \textbf{Sedative effect:} causes drowsiness and helps sleep
    \item \textbf{Muscle relaxant:} reduces muscle spasm and stiffness
    \item \textbf{Anticonvulsant:} stops or prevents seizures by calming abnormal brain activity
    \item \textbf{Amnesic effect:} reduces memory of procedures
    \item \textbf{Long half-life:} diazepam and its active metabolite remain in the body for a long time, which affects dosing and accumulation
\end{itemize}

\section*{Clinical Features}
\begin{itemize}
    \item Taken by mouth for anxiety, muscle spasm, or as an aid for procedures
    \item Given rectally or intravenously in emergencies such as prolonged seizures
    \item Effects begin within 30--60 minutes when taken by mouth
    \item Common side effects: drowsiness, dizziness, poor coordination, and impaired concentration
    \item Effects on memory and judgement, including effects on driving
    \item Withdrawal symptoms if stopped suddenly after regular use
    \item Tolerance develops with continued use, requiring higher doses for the same effect
\end{itemize}

\section*{Differential Diagnosis}
\begin{itemize}
    \item Diazepam versus other benzodiazepines, which differ in onset and duration
    \item Diazepam versus other medicines for anxiety, such as antidepressants and pregabalin
    \item Diazepam versus other antiseizure medicines
    \item Short-term versus long-term treatment strategies for anxiety and insomnia
\end{itemize}

\section*{When to Seek Medical Care}
Diazepam is a prescription medicine and should be used under medical supervision. See a doctor to discuss whether it is appropriate, how long to take it, and how to stop it safely. Never stop diazepam suddenly after regular use, as withdrawal can be dangerous.

\textbf{Call 000 (or local emergency services) for:}
\begin{itemize}
    \item Difficulty breathing, severe drowsiness, or unresponsiveness after taking diazepam (possible overdose)
    \item A seizure lasting more than five minutes, for which diazepam may be part of emergency treatment
    \item Signs of severe withdrawal, including confusion, agitation, or seizures
    \item Taking diazepam with alcohol or other sedating drugs with serious effects
\end{itemize}

\section*{Diagnosis and Investigations}
\begin{itemize}
    \item \textbf{Medical assessment:} review of the condition being treated and suitability for diazepam
    \item \textbf{Review of medicines:} checking for interactions with alcohol and other sedatives
    \item \textbf{Risk assessment:} history of dependence, falls, or substance use
    \item \textbf{Liver and kidney function:} dose adjustment may be needed in some people
    \item \textbf{Pregnancy and breastfeeding:} risks are discussed before use
\end{itemize}

\section*{Management}
\begin{itemize}
    \item \textbf{Short-term use:} diazepam is used for the shortest time needed, usually days to a few weeks
    \item \textbf{Lowest effective dose:} minimises side effects and dependence
    \item \textbf{Gradual withdrawal:} tapering under medical supervision when stopping after regular use
    \item \textbf{Alternative treatments:} psychological therapy and non-benzodiazepine medicines for anxiety
    \item \textbf{Emergency use:} rectal diazepam for seizures at home, with clear instructions from the treating team
    \item \textbf{Avoiding alcohol:} alcohol greatly increases the sedative effects of diazepam
    \item \textbf{Driving:} diazepam impairs driving, and driving is not permitted while affected
\end{itemize}

\section*{Complications}
\begin{itemize}
    \item Dependence and addiction with regular use
    \item Withdrawal symptoms, including anxiety, insomnia, and seizures on sudden stopping
    \item Falls and fractures, especially in older people
    \item Drowsiness and impaired cognition affecting daily function
    \item Overdose, especially with alcohol or other sedatives
    \item Impaired driving and workplace safety
\end{itemize}

\section*{Prognosis and Outcomes}
Diazepam is very effective for short-term relief of anxiety, muscle spasm, and seizures, and it is life-saving in status epilepticus. Used briefly and under supervision, the risks are low. Long-term use carries a significant risk of dependence, which is why prescribers limit duration. People who use diazepam long-term can be withdrawn safely with a gradual taper and support, and most manage well with alternative treatments for the underlying condition.

\section*{Prevention and Self-Care}
\begin{itemize}
    \item Take diazepam exactly as prescribed and for the shortest time needed
    \item Do not drink alcohol while taking it
    \item Do not drive or operate machinery while affected
    \item Do not share the medicine with others
    \item Seek medical advice before stopping, and taper under supervision
    \item Store diazepam safely and out of reach of children
    \item Discuss alternative treatments for anxiety with your doctor
\end{itemize}

\section*{Sources and Further Reading}
The following reputable sources provide further information for healthcare students and patients seeking reliable, current advice about diazepam.

\begin{itemize}
    \item NPS MedicineWise. \textit{Diazepam: information about benzodiazepines.} https://www.nps.org.au/
    \item Healthdirect Australia. \textit{Diazepam: uses and precautions.} https://www.healthdirect.gov.au/
    \item Therapeutic Goods Administration. \textit{Benzodiazepine safety information.} https://www.tga.gov.au/
    \item NHS. \textit{Diazepam: a benzodiazepine medicine.} https://www.nhs.uk/medicines/diazepam/
    \item MedlinePlus. \textit{Diazepam.} https://medlineplus.gov/
    \item Alcohol and Drug Foundation. \textit{Benzodiazepines: information and support.} https://adf.org.au/
\end{itemize}
